\cleardoublepage
\chapternonum{致谢}
当最后一个标点落定，抬头望见窗边油菜花正盛，金黄铺展，仿佛在无声宣告一段成长的完成。回望来路，谈不上波澜壮阔、直上青云，却也绝非一帆风顺。眼前这方书桌，见证了无数个深夜的虫鸣，被咖啡水渍晕开的草稿纸，以及在键盘敲击声中悄然生长的多肉。此刻，这些细碎的时光，仿佛化作点点星屑，散落进这一叠百余页的纸张之中。谨以此文，致谢托举我攀登学术峰峦的至亲脊梁，指引我穿越迷雾的师者明灯，以及与我共沐风雨、彼此映照的同行伙伴。

亲恩如根，深植岁月。五年前，读研将毕，我向父母提出继续攻读博士的念头，他们的回应不过一个简简单单的好字，却将所有现实的重量悄然扛在肩上。七年前我提及读硕士时如此，此前二十余年我的生涯亦是如此。他们从不多言，却始终在场。于他们而言，我仿佛一只被有意放飞的风筝，愿我飞得更高、更远，但无论飞至何处，回家时，总有温热的饭菜与整齐的被褥在等待。我亦要感谢我的妻子，在我求学之际，她为我承担起对家庭的照料，在我因研究所需却苦于缺乏设备时，她毫不犹豫拿出积蓄支持我购置所需，在我陷入困顿、思路停滞之时，她陪我走出书桌，帮我松弛心绪。她或许并不熟悉我的研究领域，却始终以她能及的倾听、建议的方式参与其中，甚至细致地协助我修改论文格式与绘图。这些看似琐碎的支持，构成了我得以继续前行的重要底气。

师者如灯，照亮幽径。初入门时，我曾困于文献迷宫，不知何为科研，亦不知何以为方向。彼时不过是一只听需求、做实验、做汇报的循环机器，进展不顺时难免有批评，写作不佳时亦常被指出不足。但也正是在一篇篇被逐句批注的稿件之上，在一次次被推翻又重建的思路之中，我逐渐学会如何辨析问题、如何推进研究。在我停滞不前时，是他帮助我理清脉络，在我基础薄弱之处，是他近乎讲授般的耐心讲解。某次同车而行，他谈及海外见闻与为人之道，我忽然意识到，导师与学生的关系，或许不止于学术，他更像是将自己人生经验折叠成一只纸船，渡我至更开阔的水域。

同僚如镜，映照年轮。从玉泉到科创，再到水博园，几处空间的辗转，桌上的摆件已静立四载春秋。它们听过无数次争论与笑谈，也见证了彼此的成长与改变。这些同行之人，使研究不再是孤立的劳作。有人在我思路枯竭时邀我暂离问题，有人在深夜与我并肩至后半夜，再一同踏着夜色归去；也有人在我困惑之时耐心倾听，即便想法尚显稚嫩，却依然给予真实的反馈。亦有同伴在论文中反复挑剔细节，使我逐渐理解何为严谨表达与分寸拿捏；有人与我结伴寻味市井烟火，让紧绷的日子得以松弛。更多名字难以一一列举，却都在不经意间构成了一段共同的时光。最好的学术共同体，从来不是孤身跋涉，而是在彼此映照中前行。
% 同僚如镜，映照成长的年轮。从玉泉到科创再到水博园，桌上的摆件已经矗立了四载春秋，他虽然听了我们四年的欢笑与争论，但给我的总是一个甜美的笑容，让我明白这些伙伴多么难得。王委员总是在我没有思路时邀请我大战一把游戏，亦或者写论文时陪我熬夜到后夜，然后喝着可乐一起回寝室。小陈学弟虽然想法稚嫩，但是总能在我困惑之时倾听我的梳理并给我意见，虽然我给他布置了许多小任务，他也总是任劳任怨地完成。还有小朱的同姓战友，亦是干饭兄弟，穿梭于各个鸡脚嘎达的地方寻觅美食。更有学姐小金和学长老杨，总爱给我论文挑刺，也是他们让我学会学术表达中的严谨和给人多留三分余地。当然还有许多许多，教我游泳的小赵，给我画饼的但毫无保留的刘哥，教我写PoC和利用漏洞的小胡，药神小桂，说话总给人不紧不慢感觉的小杨，恋爱之神小楷，爱打高尔夫的小潘，二次元之神小刘和宅男小王等等等等，说不尽道不尽都融成一股暖流成为时代的记忆。最好的学术共同体，从来不是孤军奋战，而是陪伴自己前行的镜子。

岁月如歌，亦致此刻仍不忘初心的自己。曾有人问我，临近而立之年，何以仍在求学。我给出的答案始终简单：想多学一些尚未掌握的东西。为此，我记得那些啃读晦涩书籍的夜晚，记得敲击代码至手指僵硬的时刻，也记得一次次演练汇报的反复推敲。那些写满思路的草稿、记录实验的数据文档，以及被嫌弃不够美观的幻灯片，最终都沉淀为我生命的一部分。它们让我明白，学术之路从无坦途，也正因其崎岖，登高之后的风景才更为绚烂。所谓博士，不过是广阔世界中的一个微小切片，更像一个带着好奇心不断发问的婴儿，因为真正的社会人生还未开始。愿此去仍能保有赤子之心，在未知中前行，在知识的原野上，继续做一个追光的赶路人。
% 岁月如歌，致不忘初心的自己。此刻回望，清楚地记得有人曾问我，为什么快 30 了还在读书，我给的回答永远是，想多学点不会的东西。为了这个初心，我想起多少个日夜啃得砖头黑皮书，在电脑桌前打得手指变形的代码，在演说之前反复演练和背稿，那写着画着各种想法的草稿，记录着各种知识点和实验的数据和文档，以及做的总被人嫌弃丑的 PPT，都成为了我生命力最重要的印记。这都让我明白，学术之路，从无坦途，正是山路崎岖，才能在攀登之后看到最好的风景。所谓博士，不过是万千世界中一个微不足道的缩影，更是一个带着好奇的新生儿，对这个世界提问者，到底为什么是这样。纸短情长，墨浅意深，愿我能永葆赤子之心，继续抱有好奇，多学点东西，在知识的原野上做一个永远追光的赶路人。
